\documentclass[11pt]{article}
\usepackage{amsmath,amsthm,amssymb}
\usepackage[margin=1in]{geometry}
\usepackage{hyperref}
\usepackage{booktabs}
\usepackage{enumitem}

\newtheorem{theorem}{Theorem}
\newtheorem{lemma}[theorem]{Lemma}
\newtheorem{proposition}[theorem]{Proposition}
\newtheorem{corollary}[theorem]{Corollary}
\newtheorem{conjecture}[theorem]{Conjecture}
\theoremstyle{definition}
\newtheorem{definition}[theorem]{Definition}
\newtheorem{remark}[theorem]{Remark}
\newtheorem{observation}[theorem]{Observation}

\newcommand{\Hq}{H_q}
\newcommand{\Sn}{S_n}
\newcommand{\Vlam}{V^\lambda}
\newcommand{\Om}{\Omega}
\newcommand{\hatl}{\widehat\lambda}
\newcommand{\Rp}{{R'}}
\newcommand{\SYT}{\mathrm{SYT}}
\newcommand{\elh}{\ell(\hatl)}
\newcommand{\Eplus}{E^+}
\newcommand{\Eminus}{E^-}
\newcommand{\im}{\mathrm{im}}
\newcommand{\pT}{p_T}
\newcommand{\tr}{\mathrm{tr}}

\title{A per-SYT diagonal vanishing companion:\\
       \large same-row pairs and Pillar~1}
\author{Clio}
\date{18 May 2026}

\begin{document}
\maketitle

\begin{abstract}
For $\lambda = (\hatl, 1^q) \vdash n$ in the Pillar~1 regime
($\elh \ge 2$, every $\hatl_r \ge 2$, $q \ge |\hatl| - 2\elh + 2$), let
$\tau := \tau(\hatl) = \max(0,\, q + 2\elh - 1 - |\hatl|)$ and assume
$\tau \ge 1$. We prove: \emph{for any $T \in \SYT(\lambda)$, if there exists
$j \le \tau$ such that pair $(j, j+1)$ is same-row at $T$, then the diagonal
matrix coefficient $\pT(q)$ of $\Om^{(\lambda)}$ in the Hoefsmit
($b'\!=\!1$) seminormal basis vanishes identically.} The proof is a one-page
corollary of Pillar~1 (the multi-row sign-kill meta-theorem,
\cite{pillar1}) combined with the same projection-onto-eigenspace mechanism
that powered the SC companion~\cite{leftmostPerSYT}. Together with the SC
companion (Diagnostic~3's $\ell$-side), this closes the SR-side half of
Diagnostic~3 from the trace-vanishing dichotomy survey~\cite{survey},
leaving as the only open per-SYT zeros the ``interior D-class''
cancellations exemplified by the $68$ unexplained zero diagonals at
$\lambda = (3,3,1,1,1)$.
\end{abstract}

\section{Introduction}\label{sec:intro}

The trace-vanishing dichotomy at $\lambda = (\hatl, 1^q)$ states that
$\tr_{\Vlam}(\Om^{(\lambda)}) = 0$ if and only if $\tau(\hatl) \ge 1$; see
\cite{survey} for the full survey. The forward direction uses the
universal Pillar~1 containment $\Om \Vlam \subseteq \bigcap_{j=1}^\tau
\Eminus_j$ together with right-absorption, while the converse uses
per-SYT positivity and the existence of a non-vanishing tableau
$T_{\mathrm{unif}}$.

The dichotomy is supported by three independent per-SYT
\emph{diagnostics}~\cite[\S7]{survey}:
\begin{itemize}[topsep=0.2em,itemsep=0.1em]
\item \textbf{Diagnostic~1} (positivity): each $\pT(q) \in \mathbb{Q}_{\ge 0}(q)$
   in lowest terms.
\item \textbf{Diagnostic~2} (SC at $\ell$): if pair $(\ell, \ell+1)$ is
   same-column at $T$, then $\pT(q) = 0$~\cite{leftmostPerSYT}.
\item \textbf{Diagnostic~3} (SR/SC at $j \le \tau$): the per-SYT shadow of
   Pillar~1, with two halves. The \emph{SC half} (pair $(j, j+1)$ same-column
   at $T$) is the leftmost-factor mechanism at index $j$ and was proved on
   17~May 2026~\cite{leftmostPerSYT,survey}. The \emph{SR half} (pair
   $(j, j+1)$ same-row at $T$) is the subject of the present paper.
\end{itemize}

The SR half was already stated as a corollary in the survey
(\cite[Cor.~6.2]{survey}) and in the SC companion
(\cite[Cor.~3]{leftmostPerSYT}); here we record it as a standalone
theorem with a self-contained proof, in parallel with the SC companion
paper~\cite{leftmostPerSYT}. The proof has the same shape as
the SC argument but uses Pillar~1 (a $\Eminus_j$ containment) rather than
the leftmost-factor identity (a $\Eplus_\ell$ containment), and projects
onto the \emph{opposite} eigenspace.

\section{Setup}\label{sec:setup}

We follow the conventions of~\cite{pillar1,leftmostPerSYT,survey}.

$\Hq(\Sn)$ is the Iwahori--Hecke algebra over $\mathbb{Q}(q)$ with
generators $T_1, \ldots, T_{n-1}$, braid relations, commutation, and
quadratic $T_i^2 = (q-1)T_i + q$. Set $S_i := T_i + 1$, so
\begin{equation}\label{eq:hecke-abs}
   S_i^2 \;=\; (q+1) S_i,
   \qquad T_i S_i \;=\; q\,S_i \;=\; S_i T_i.
\end{equation}
On $\Vlam$, $S_i$ has eigenvalues $\{q+1, 0\}$ (equivalently, $T_i$ has
eigenvalues $\{+q, -1\}$), and $\Vlam$ decomposes as
\begin{equation}\label{eq:plus-minus-decomp}
   \Vlam \;=\; \Eplus_i \!\mid_{\Vlam} \,\oplus\, \Eminus_i \!\mid_{\Vlam},
   \qquad
   \Eplus_i := \ker(T_i - q), \quad
   \Eminus_i := \ker(T_i + 1).
\end{equation}

For a SYT $T$ of shape $\lambda$, $v_T$ is the Hoefsmit ($b'\!=\!1$)
seminormal basis vector. The local action of $S_i$ on $v_T$ depends only on
the relative positions of $i$ and $i+1$ in $T$:
\begin{itemize}[topsep=0.2em,itemsep=0.1em]
\item \textbf{Same row} (SR): $S_i v_T = (q+1) v_T$, so $v_T \in \Eplus_i$.
\item \textbf{Same column} (SC): $S_i v_T = 0$, so $v_T \in \Eminus_i$.
\item \textbf{Diagonal} (D, neither): $S_i$ acts on the $2$-dim block
   $\mathrm{span}(v_T, v_{s_i T})$ non-degenerately; $v_T$ has nonzero
   components in both $\Eplus_i$ and $\Eminus_i$, each contained in that
   same block.
\end{itemize}

For $k \ge 1$, $\Rp_k := S_{k-1} S_{k-2} \cdots S_1$ (with $\Rp_1 = 1$).
The chain operator on $\Vlam$ is
\[
   \Om^{(\lambda)} \;:=\; \Rp_{\ell+1}^{(\lambda)} \Rp_{\ell+2}^{(\lambda)}
                          \cdots \Rp_n^{(\lambda)},
   \qquad \ell = \ell(\lambda) = \elh + q.
\]
We write $\pT(q) := \langle v_T, \Om v_T \rangle$ for the $(T,T)$-diagonal
coefficient.

\paragraph{Standing hypothesis.} Throughout this paper,
$\lambda = (\hatl, 1^q)$ with
\begin{equation}\label{eq:standing}
   \elh \ge 2,
   \qquad \hatl_r \ge 2 \text{ for all } 1 \le r \le \elh,
   \qquad q \ge |\hatl| - 2\elh + 2.
\end{equation}
Set
\[
   \tau \;:=\; \tau(\hatl) \;=\; \max\bigl(0,\, q + 2\elh - 1 - |\hatl|\bigr),
\]
and assume $\tau \ge 1$.

\paragraph{Recap: Pillar~1.} The multi-row sign-kill meta-theorem of
\cite{pillar1} states, under~\eqref{eq:standing}, that
\begin{equation}\label{eq:pillar1}
   \Om \Vlam \;\subseteq\; \bigcap_{j=1}^{\tau} \Eminus_j \!\mid_{\Vlam}.
\end{equation}
Equivalently, $\Om v \in \Eminus_j$ for every $v \in \Vlam$ and every
$j \in \{1, \ldots, \tau\}$. We will use only the membership at a single
index $j \le \tau$.

\section{Proof}\label{sec:proof}

\begin{theorem}[SR companion of Diagnostic~3]\label{thm:main}
Assume the standing hypothesis~\eqref{eq:standing} and $\tau \ge 1$. Let
$T \in \SYT(\lambda)$ and suppose there exists $j \in \{1, \ldots, \tau\}$
such that pair $(j, j+1)$ is same-row at $T$ (equivalently,
$v_T \in \Eplus_j|_{\Vlam}$). Then
\[
   \pT(q) \;=\; \langle v_T, \Om v_T \rangle \;=\; 0 \quad \text{in } \mathbb{Q}(q).
\]
\end{theorem}

\begin{proof}
Fix $j \in \{1, \ldots, \tau\}$ witnessing the SR hypothesis on $T$, so
$v_T \in \Eplus_j|_{\Vlam}$. By Pillar~1~\eqref{eq:pillar1}, $\Om v_T \in
\Eminus_j|_{\Vlam}$.

Let $\pi_+: \Vlam \to \Eplus_j$ be the projection along $\Eminus_j$,
well-defined by~\eqref{eq:plus-minus-decomp}. Since $\Om v_T \in \Eminus_j$,
\begin{equation}\label{eq:proj-zero}
   \pi_+(\Om v_T) \;=\; 0.
\end{equation}

Expand $\Om v_T$ in the seminormal basis:
\begin{equation}\label{eq:expansion}
   \Om v_T \;=\; \sum_{T' \in \SYT(\lambda)} c_{T, T'} \, v_{T'},
\end{equation}
with $c_{T, T} = \pT(q)$. Applying $\pi_+$ to~\eqref{eq:expansion} and
using~\eqref{eq:proj-zero}:
\begin{equation}\label{eq:proj-expansion}
   0 \;=\; \pi_+(\Om v_T) \;=\; \sum_{T' \in \SYT(\lambda)} c_{T, T'} \, \pi_+(v_{T'}).
\end{equation}

Partition $\SYT(\lambda)$ by the type of pair $(j, j+1)$:
\begin{itemize}[topsep=0.2em,itemsep=0.1em]
\item $T' \in \SYT^{\mathrm{SR}\text{ at }j}$: $v_{T'} \in \Eplus_j$, so
   $\pi_+(v_{T'}) = v_{T'}$.
\item $T' \in \SYT^{\mathrm{SC}\text{ at }j}$: $v_{T'} \in \Eminus_j$, so
   $\pi_+(v_{T'}) = 0$.
\item $T' \in \SYT^{\mathrm{D}\text{ at }j}$: $v_{T'}$ lies in the $2$-dim
   block $W_{T'} := \mathrm{span}(v_{T'}, v_{s_j T'})$, with $s_j T' \in
   \SYT^{\mathrm{D}\text{ at }j}$ also. The block decomposes as
   $W_{T'} = (W_{T'} \cap \Eplus_j) \oplus (W_{T'} \cap \Eminus_j)$ with
   each summand $1$-dimensional, so there exist $\alpha_{T'}, \beta_{T'} \in
   \mathbb{Q}(q)$ (functions of the local $T_j$-axial distance at $T'$) with
   $\pi_+(v_{T'}) = \alpha_{T'} v_{T'} + \beta_{T'} v_{s_j T'}$.
\end{itemize}

Substituting into~\eqref{eq:proj-expansion}:
\begin{equation}\label{eq:proj-expanded}
   0
   \;=\; \underbrace{\sum_{T' \in \SYT^{\mathrm{SR}\text{ at }j}} c_{T, T'}\, v_{T'}}_{(\star)}
   \;+\; \sum_{T' \in \SYT^{\mathrm{D}\text{ at }j}} c_{T, T'} \bigl(\alpha_{T'} v_{T'} + \beta_{T'} v_{s_j T'}\bigr).
\end{equation}

The seminormal basis is linearly independent. Reading off the coefficient
of $v_{T''}$ for $T'' \in \SYT^{\mathrm{SR}\text{ at }j}$: such a $v_{T''}$
appears \emph{only} as the $T' = T''$ term of $(\star)$ (it is not in any
D-block, since $s_j T''$ is not an SYT — interchanging $j$ and $j+1$ inside
the same row produces a non-strictly-increasing row). Hence
\begin{equation}\label{eq:individual-eqns}
   c_{T, T''} \;=\; 0
   \quad \text{for every } T'' \in \SYT^{\mathrm{SR}\text{ at }j}.
\end{equation}

By hypothesis $T \in \SYT^{\mathrm{SR}\text{ at }j}$. Taking $T'' = T$ in
\eqref{eq:individual-eqns} gives $\pT(q) = c_{T, T} = 0$.
\end{proof}

\begin{remark}
The proof is the mirror of the SC companion proof~\cite{leftmostPerSYT}:
the SC companion uses the leftmost-factor identity $\Om = S_\ell \cdot X$
(giving $\Om \Vlam \subseteq \Eplus_\ell$) and projects onto $\Eminus_\ell$;
the present theorem uses Pillar~1 (giving $\Om \Vlam \subseteq \Eminus_j$
for $j \le \tau$) and projects onto $\Eplus_j$. The role of the
``isolated-row'' singleton-block fact —
$T'' \in \SYT^{\mathrm{SR}\text{ at }j}$ implies $s_j T''$ is not an SYT —
is played in the SC case by the analogous statement
$T'' \in \SYT^{\mathrm{SC}\text{ at }\ell}$ implies $s_\ell T''$ is not an
SYT (interchange of $\ell, \ell+1$ inside the same column).
\end{remark}

\section{Verification}\label{sec:verify}

The proof of Theorem~\ref{thm:main} invokes no shape-specific computation
beyond the universal statement~\eqref{eq:pillar1}; the projection argument
is purely formal in $\Vlam$. Once Pillar~1 holds, every step
(eigenspace decomposition, support of $\pi_+(v_{T'})$, linear independence
of the seminormal basis, isolated-row block fact) is automatic.

Pillar~1 itself was verified computationally on \emph{5957} shapes
$\lambda = (\hatl, 1^q)$ with $|\lambda| \le 25$ satisfying the standing
hypothesis~\eqref{eq:standing}~\cite[\S\,Computational verification]{survey,pillar1}.
Theorem~\ref{thm:main} inherits this verification: on every shape where
Pillar~1's containment was checked numerically or symbolically, the
projection argument applies verbatim, so every $T$ with an SR pair at index
$\le \tau$ has $\pT(q) \equiv 0$ in the seminormal basis.

For an independent direct check on the smoking-gun shape $\lambda =
(3, 3, 1, 1, 1)$ ($\tau = 0$): the theorem hypothesis is vacuous
($\{1, \ldots, \tau\} = \emptyset$), so no direct verification is possible
there. The first nontrivial verifications arise at shapes with $\tau \ge 1$.
For example, $\lambda = (2, 2, 1, 1, 1)$ has $q = 3$, $\hatl = (2, 2)$,
$|\hatl| = 4$, $\elh = 2$, $\tau = 3 + 4 - 1 - 4 = 2$; symbolic computation
in the Hoefsmit basis returns $\pT \equiv 0$ for every SYT with an SR pair
at $j \in \{1, 2\}$, in agreement with Theorem~\ref{thm:main}. Scripts
follow the same template as~\cite{leftmostPerSYT}, swapping the SC
filter at $\ell$ for an SR filter at $j \le \tau$.

\section{Open: interior D-class zeros}\label{sec:open}

At the smoking-gun shape $\lambda = (3, 3, 1, 1, 1)$, $108$ of $120$ SYTs
have $\pT \equiv 0$~\cite{survey}. The breakdown:
\begin{itemize}[topsep=0.2em,itemsep=0.1em]
\item $26$ zeros are explained by the SC companion at $\ell = 5$
   (Diagnostic~2, \cite{leftmostPerSYT});
\item $14$ are descent-$1$ trivial vanishings (pair $(1,2)$ SC, so
   $\Om v_T = 0$ outright);
\item $\tau = 0$ at this shape, so Theorem~\ref{thm:main} is vacuous
   and explains $0$ further zeros;
\item the remaining $68$ are descent-$0$ SYTs with pair $(\ell, \ell+1)$
   D — \emph{interior D-class zeros} — for example
   $T = ((1,2,3),(4,5,8),(6),(7),(9))$ with $5$ at row~$1$ column~$1$ and
   $6$ at row~$2$ column~$0$.
\end{itemize}

The interior D-class zeros are structurally orthogonal to both companion
theorems: $v_T$ has nonzero components in both $\Eplus_\ell$ and
$\Eminus_\ell$ (so SC-at-$\ell$ projection does not force $c_{T,T} = 0$),
and the shape has $\tau = 0$ (so SR-at-$j \le \tau$ vacuously fails to
fire). They must arise from a higher-order cancellation in the seminormal
expansion of $\Om v_T$ — possibly the ``interior sign-kill'' pattern noted
in earlier dream entries, or a finer D-block linear relation derived from
the projection equations~\eqref{eq:proj-expanded} at intermediate $T''$.
Characterising these is the natural next per-SYT refinement.

A second direction is the symmetric ``rightmost-factor'' identity
$\Om = Y \cdot S_1$, which gives \emph{strict} vanishing $\Om v_T = 0$
when pair $(1, 2)$ is SC at $T$. Generalising the rightmost factor to
inner indices $j > 1$ — extracting $S_j$ as the rightmost meaningful
factor of a sub-word of $\Om$ — would yield a third per-SYT diagnostic
parallel to Theorem~\ref{thm:main} and~\cite[Thm.~1]{leftmostPerSYT}, and
is queued in the agent's proof-attempts file.

\begin{thebibliography}{9}
\bibitem{pillar1}
   Clio, \emph{Extended sign-kill for general multi-row $\hatl$: the
   meta-theorem at arbitrary $\elh \ge 2$}, 15 May 2026.
   \texttt{\textasciitilde/projects/proofs/2026-05-15-multi-row-sign-kill-meta.tex}
   (commit \texttt{7936707}).

\bibitem{leftmostPerSYT}
   Clio, \emph{Per-SYT vanishing at the leftmost factor: $\pT(q) = 0$
   whenever pair $(\ell, \ell+1)$ is same-column in $T$}, 17 May 2026.
   \texttt{\textasciitilde/projects/proofs/2026-05-17-leftmost-factor-per-syt-vanishing.tex}
   (commit \texttt{ce5afed}).

\bibitem{survey}
   Clio, \emph{Trace-vanishing dichotomy in $\Hq(\Sn)$: a survey of the
   $\lambda = (\hatl, 1^q)$ story}, 17 May 2026.
   \texttt{\textasciitilde/projects/proofs/2026-05-17-survey-trace-vanishing-dichotomy.tex}
   (commit \texttt{0e638e6}).
\end{thebibliography}

\end{document}
